\documentclass[12pt]{report}
\usepackage[brazilian]{babel}
\usepackage{csquotes}
\usepackage{makeidx}
\usepackage{glossaries}
\usepackage[
  perfil=academico,
  tipo-trabalho=tese,
  impressao=frente-e-verso,
  citacoes=autor-data
]{abntex3}

\makeindex
\makeglossaries
\newglossaryentry{preservacao}{
  name={preservação digital},
  description={processos destinados a manter documentos digitais acessíveis}
}
\addbibresource{referencias-exemplo.bib}

\ABNTEXsetup{
  titulo={Modelo canônico completo de tese},
  subtitulo={base auditável com todos os elementos implementados},
  autoria={Maria da Silva},
  instituicao={Universidade Exemplo},
  natureza={Tese},
  objetivo={apresentada para obtenção do título de Doutora},
  area={Ciência da Informação},
  orientacao={Prof. Dr. João Souza},
  coorientacao={Profa. Dra. Ana Santos},
  volume={1},
  local={Recife},
  data={2026}
}

\ABNTEXabreviatura{ed.}{edição}
\ABNTEXsigla{ABNT}{Associação Brasileira de Normas Técnicas}
\ABNTEXsigla{PDF}{Portable Document Format}
\ABNTEXsimbolo{$n$}{Quantidade de documentos}
\ABNTEXsimbolo{$\sum$}{Somatório}

\begin{document}

% PARTE EXTERNA
\ABNTEXacademicoexterna
\ABNTEXlombada[Tese 2026][Universidade Exemplo]

% ELEMENTOS PRÉ-TEXTUAIS
\ABNTEXacademicopretextual
\ABNTEXfolhaderosto
\ABNTEXfichacatalografica{%
  Ponto de extensão para a ficha preparada por profissional responsável.
}
\ABNTEXerrata
  {SILVA, Maria da. \textbf{Modelo canônico completo de tese}. 2026.
   Tese — Universidade Exemplo, Recife, 2026.}
  {\begin{center}
     \begin{tabular}{llll}
       Folha & Linha & Onde se lê & Leia-se\\
       10 & 4 & preservacão & preservação
     \end{tabular}
   \end{center}}
\ABNTEXfolhadeaprovacao
  {30 de julho de 2026}
  {%
    \ABNTEXmembrobanca
      {Prof. Dr. João Souza}{Orientador}{Universidade Exemplo}
    \ABNTEXmembrobanca
      {Profa. Dra. Ana Santos}{Coorientadora}{Universidade Exemplo}
    \ABNTEXmembrobanca
      {Prof. Dr. Carlos Lima}{Examinador externo}{Instituto Exemplo}
  }
\ABNTEXdedicatoria{Às pessoas que preservam e compartilham conhecimento.}
\ABNTEXagradecimentos{%
  À comunidade de software livre e às pessoas responsáveis pela avaliação
  normativa independente deste modelo.
}
\ABNTEXepigrafe[Autoria demonstrativa]{%
  Um documento auditável explicita tanto seu conteúdo quanto suas escolhas.
}
\ABNTEXresumo
  {Este modelo demonstra todos os elementos públicos aplicáveis a uma tese,
   incluindo listas, divisões textuais, referências, glossário, apêndice,
   anexo e índice.}
  {tese, normalização, auditoria, LaTeX}
\ABNTEXresumoemlinguaestrangeira
  {This template demonstrates every public element applicable to a thesis,
   including lists, textual divisions, references, glossary, appendix,
   annex, and index.}
  {thesis, standardization, audit, LaTeX}
\ABNTEXlistadeilustracoes
\ABNTEXlistadetabelas
\ABNTEXlistadeabreviaturasesiglas
\ABNTEXlistadesimbolos
\ABNTEXsumario

% ELEMENTOS TEXTUAIS
\ABNTEXacademicotextual
\chapter{Introdução}
\index{auditoria}
Este capítulo delimita o tema, apresenta problema, hipótese, objetivos e
justificativa. A \gls{preservacao} fornece o domínio de demonstração.
Uma citação indireta usa a integração bibliográfica
\parencite{associacao2025}.

\section{Problema de pesquisa}
Como tornar modelos documentais simultaneamente reutilizáveis e auditáveis?

\section{Hipótese}
A explicitação dos elementos e a compilação automatizada facilitam a revisão.

\section{Objetivos}
O objetivo geral é demonstrar o perfil. Os objetivos específicos são
inventariar elementos, exercitar a API e validar a distribuição.

\chapter{Referencial teórico}
\section{Normalização documental}
O estilo autor-data também permite citação integrada:
\textcite{associacao2025}.
\begin{ABNTEXcitacaolonga}
Este texto meramente demonstrativo ocupa o ambiente destinado a uma citação
direta longa. A fonte real deve ser indicada pela autoria do novo documento.
\end{ABNTEXcitacaolonga}
\ABNTEXnota{Nota demonstrativa composta pela API pública do layout.}

\chapter{Metodologia}
\section{Delineamento}
O estudo demonstrativo utiliza revisão documental e ensaios de compilação.
\section{Procedimentos}
Cada elemento público é composto e posteriormente verificado.

\chapter{Resultados e discussão}
\section{Resultados}
\begin{figure}[ht]
  \centering
  \rule{0.6\linewidth}{2cm}
  \caption{Ilustração demonstrativa do fluxo}
  \ABNTEXfonte{elaborada pela autora.}
\end{figure}
\begin{table}[ht]
  \centering
  \caption{Cobertura dos elementos}
  \begin{tabular}{ll}
    Grupo & Estado\\
    Pré-textuais & demonstrados\\
    Pós-textuais & demonstrados
  \end{tabular}
  \ABNTEXfonte{elaborada pela autora.}
\end{table}
\section{Discussão}
Os resultados devem ser confrontados com o referencial adotado.

\chapter{Conclusão}
Retomam-se objetivos e resultados, registrando limitações e trabalhos futuros.

% ELEMENTOS PÓS-TEXTUAIS
\ABNTEXacademicopostextual
\ABNTEXbibliografia
\ABNTEXglossario
\ABNTEXappendix{Instrumento elaborado pela autora}
Conteúdo de responsabilidade da autoria do trabalho.
\ABNTEXannex{Documento externo reproduzido para demonstração}
Conteúdo não elaborado pela autoria do trabalho.
\ABNTEXindice[Índice de assuntos]

\end{document}
